\chapter{The Pipeline, Line by Line}
\label{ch:pipeline}

Everything in this book converges on one 70-line Groovy file. A
declarative Jenkinsfile lives in the repository it builds --- pipeline
as code, versioned and reviewed like everything else.

\section{Preamble and environment}

\begin{lstlisting}
pipeline {
  agent any                                  (*@\co{1}@*)
  options {
    timestamps()
    disableConcurrentBuilds()                (*@\co{2}@*)
    buildDiscarder(logRotator(numToKeepStr: '10'))
  }
  triggers { pollSCM('H/5 * * * *') }        (*@\co{3}@*)
  environment {
    JAVA_HOME = "${env.JDK25_HOME}"          (*@\co{4}@*)
    PATH      = "${env.JDK25_HOME}/bin:${env.PATH}"
    TAG       = "${env.GIT_COMMIT.take(7)}"  (*@\co{5}@*)
    NS        = 'ts'
    REGISTRY  = 'localhost:5000'
    SERVICES  = 'config-server discovery-server course-service api-gateway auth-service dictionary-service notification-service'
    MODULES   = 'common,config-server,discovery-server,course-service,api-gateway,auth-service,dictionary-service,notification-service'
  }
\end{lstlisting}

\begin{description}
  \item[\co{1}] Any executor on the controller --- Chapter~22's
  homelab compromise in one word.
  \item[\co{2}] Two builds at once on two cores would thrash; queued
  runs also collapse (a new commit during a build supersedes waiting
  ones).
  \item[\co{3}] Polling because NAT blocks webhooks. The \code{H} is a
  hash-spread: ``some fixed minute in each 5,'' so many jobs don't all
  poll at :00.
  \item[\co{4}] The per-pipeline JDK switch the controller image
  prepared: builds see Java\,25, Jenkins itself stays on 21.
  \item[\co{5}] The deploy tag is the short git SHA --- every image
  traceable to its commit, every deploy pinned to an immutable tag
  (Chapter~13's tags-vs-digests lesson applied).
  \item[\code{SERVICES}/\code{MODULES}] the pipeline's registry of what
  exists: \code{MODULES} feeds Maven (comma-separated, includes
  \code{common} which is built but produces no image);
  \code{SERVICES} feeds kubectl (space-separated, deployables only).
  Adding a service to the cluster means extending both --- the one
  manual step per new service.
\end{description}

\section{The three stages}

\begin{lstlisting}
stage('Build') {
  steps {
    dir('Project/TeacherSupporter') {          (*@\co{1}@*)
      sh 'mvn -B -DskipTests -pl ${MODULES} -am verify'  (*@\co{2}@*)
    } } }
stage('Push images') {
  steps {
    dir('Project/TeacherSupporter') {
      sh 'mvn -B -DskipTests -pl ${MODULES} -am package jib:build -Djib.to.tags=${TAG}'  (*@\co{3}@*)
    } } }
stage('Deploy') {
  steps {
    dir('Project/TeacherSupporter') {
      withCredentials([file(credentialsId: 'kubeconfig-ts',
                            variable: 'KUBECONFIG')]) {   (*@\co{4}@*)
        sh '''
          set -eu                                         (*@\co{5}@*)
          kubectl apply -k deploy/k8s/base                (*@\co{6}@*)
          for s in ${SERVICES}; do
            kubectl -n ${NS} set image deployment/${s} \
                ${s}=${REGISTRY}/${s}:${TAG}              (*@\co{7}@*)
          done
          for s in ${SERVICES}; do
            kubectl -n ${NS} rollout status \
                deployment/${s} --timeout=300s            (*@\co{8}@*)
          done
        '''
      } } } }
\end{lstlisting}

\begin{description}
  \item[\co{1}] The repository root is a study monorepo; the project
  lives in a subdirectory, so every stage \code{dir()}s into it.
  \item[\co{2}] \code{-pl} (project list) + \code{-am} (also make
  dependencies): build only the listed modules and what they need ---
  monorepo economics on two cores. \code{-B} is batch mode (no
  interactive output). \code{-DskipTests} is the consciously accepted
  gate-weakening from Chapter~21's pitfall box.
  \item[\co{3}] One command builds and pushes all seven images ---
  Chapter~13 in a single line. \code{-Djib.to.tags} adds the SHA tag
  alongside the SNAPSHOT tag. The \code{package} in front of
  \code{jib:build} is not decoration --- see the pitfall below; it
  cost two failed builds on the day milestones 5 and 6 shipped.
  \item[\co{4}] The file credential materializes as a temp file;
  \code{KUBECONFIG} is the env var kubectl reads. Outside this block
  the file does not exist; in the log its content never appears.
  \item[\co{5}] Fail on any error and on unset variables --- without
  it, a failed apply would be sailed past and ``deployed'' reported.
  \item[\co{6}] Converge the cluster on the manifests --- idempotent
  (Chapter~14), so it is safely re-run every build whether or not
  anything changed.
  \item[\co{7}] Repoint each Deployment at the SHA-tagged image. This
  is what actually triggers rollouts --- and note the manifests
  themselves keep the SNAPSHOT tag, so \emph{apply} alone would not
  redeploy; the pipeline's pinning and the manifests' default coexist.
  \item[\co{8}] The verification gate: block until each rollout
  completes; fail the build if any pod never goes Ready within five
  minutes. This line is what makes the pipeline's green mean
  ``running,'' not ``commands executed'' --- it delegates truth to the
  probes of Chapter~17.
\end{description}

\begin{pitfall}[Could not find artifact com.ts:common:jar:1.0.0-SNAPSHOT]
Symptom: the Build stage passes, then the Push stage fails on the
first module that depends on \code{common}, unable to resolve a jar
that the previous stage just built. Cause: \code{mvn jib:build} is a
\emph{bare goal} --- Maven runs it directly with no lifecycle phases,
so nothing in the reactor is compiled or packaged in that session, and
sibling dependencies fall back to the local repository, where
\code{common} was never installed (\code{verify} stops short of
\code{install}). Modules without a sibling dependency (config-server,
discovery-server) sail through, which makes the failure look
module-specific. Fix: put a phase in front --- \code{package jib:build}
--- so the reactor packages \code{common} before Jib asks for it.
General rule: a plugin goal that needs other modules' artifacts must
ride a lifecycle phase, not be invoked alone.
\end{pitfall}

Two loops, not one, is a small correctness point: all rollouts start
first, then all are awaited --- otherwise service N's wait serializes
service N+1's start, stretching total deploy time for no gain.

\section{Failure handling}

\begin{lstlisting}
post {
  success { echo "Deployed ${TAG} to namespace ${NS}." }
  failure {
    dir('Project/TeacherSupporter') {
      withCredentials([...]) {
        sh 'kubectl -n ${NS} get pods || true'
      } } } }
\end{lstlisting}

On failure, the build log ends with a pod listing --- the first command
a human would run, already run. The \code{|| true} keeps the
diagnostic itself from masking the build's real exit status.

\section{The same pipeline, elsewhere}

For calibration, the identical shape as GitHub Actions:

\begin{lstlisting}[language=yaml]
on: { push: { branches: [main] } }        # webhook, not polling
jobs:
  deploy:
    runs-on: ubuntu-latest                # disposable agent, not controller
    steps:
      - uses: actions/checkout@v4
      - uses: actions/setup-java@v4
        with: { distribution: temurin, java-version: 25 }
      - run: mvn -B -pl $MODULES -am verify
      - run: mvn -B -pl $MODULES -am jib:build -Djib.to.tags=${GITHUB_SHA::7}
      - run: kubectl apply -k deploy/k8s/base && ...
        env: { KUBECONFIG_DATA: ${{ secrets.KUBECONFIG }} }
\end{lstlisting}

Same stages, same tag scheme, same secret-binding idea. The
differences are the infrastructure trade: GitHub's runners are
disposable VMs (agents done right, for free) but they live on the
internet --- so this variant would need the registry and cluster
reachable from outside, the exact constraint that shaped the homelab
design. Pipelines are portable; \emph{network topology is what actually
differs between shops}.

\begin{quiz}
\begin{enumerate}
  \item Why tag deploys with the git SHA when the SNAPSHOT tag also
  exists? Name the failure mode SNAPSHOT-only deploys invite
  (Chapter~13) and what traceability the SHA adds.
  \item Explain \code{-pl ... -am} and what it saves in a monorepo.
  What breaks if you forget to extend \code{MODULES} for a new service?
  \item Which single line makes this pipeline's success mean ``pods
  are Ready,'' and which chapter's machinery does it lean on?
  \item Why two loops in the deploy stage instead of
  set-image-then-wait per service?
  \item List three things that change and three that do not when this
  Jenkinsfile is translated to GitHub Actions.
\end{enumerate}
\end{quiz}
